\section{Evaluation and Statistics}

The evaluation layer converts LAMMPS log summaries into adsorption isotherm tables, reports, and plots.
The simulation molecule count is interpreted as absolute adsorption, and excess adsorption can be calculated when pore volume and gas density are available.

\subsection{Absolute Adsorption}

The absolute loading is calculated from the mean number of adsorbate molecules in the simulation cell:

\begin{equation}
    n_{\mathrm{abs}} =
    \frac{1000 \cdot \langle N_{\mathrm{ads}}\rangle}{m_{\mathrm{framework}}}
\end{equation}

where:

\begin{itemize}
    \item $n_{\mathrm{abs}}$ is the absolute adsorption loading in mol kg$^{-1}$,
    \item $\langle N_{\mathrm{ads}}\rangle$ is the mean number of adsorbate molecules per simulation cell,
    \item $m_{\mathrm{framework}}$ is the framework mass in atomic mass units.
\end{itemize}

The numerical values of mol kg$^{-1}$ and mmol g$^{-1}$ are identical.

\subsection{Excess Adsorption}

Excess adsorption is calculated as:

\begin{equation}
    n_{\mathrm{ex}} =
    n_{\mathrm{abs}} - \rho_{\mathrm{gas}} V_{\mathrm{pore}}
\end{equation}

where:

\begin{itemize}
    \item $n_{\mathrm{ex}}$ is the excess adsorption loading in mol kg$^{-1}$,
    \item $n_{\mathrm{abs}}$ is the absolute adsorption loading in mol kg$^{-1}$,
    \item $\rho_{\mathrm{gas}}$ is the bulk gas molar density in mol m$^{-3}$,
    \item $V_{\mathrm{pore}}$ is the pore volume in m$^3$ kg$^{-1}$.
\end{itemize}

For IRMOF-1, the current automatic pore volume source is:

\begin{lstlisting}
CRAFTED-2.0.0/RAC_DBSCAN/CRAFTED_MOF_geometric.csv:AV_cm^3/g
\end{lstlisting}

with:

\begin{equation}
    V_{\mathrm{pore}} = 0.845545~\mathrm{cm^3~g^{-1}}
    = 0.000845545~\mathrm{m^3~kg^{-1}}.
\end{equation}

The gas density is calculated with CoolProp using the configured backend.
For production evaluation, the current backend is \texttt{HEOS}.

\subsection{Current Evaluation Outputs}

For each evaluated run, the pipeline writes:

\begin{itemize}
    \item \texttt{evaluated\_isotherm.csv},
    \item \texttt{evaluation\_report.md},
    \item \texttt{isotherm\_comparison.png}.
\end{itemize}

For comparison against all available references, the pipeline writes:

\begin{itemize}
    \item \texttt{evaluations/reference\_candidates.json},
    \item \texttt{evaluations/combined\_reference\_comparison.csv},
    \item \texttt{evaluations/combined\_reference\_comparison.png},
    \item one evaluation subdirectory per active reference.
\end{itemize}

\subsection{Independent Replicates and Convergence}

Independent GCMC replicates are configured with \texttt{simulation.seeds}.
Every seed produces separate input, log, dump, and summary files for every pressure point.
The pipeline aggregates the replicate means and reports the sample standard deviation, standard error, and a two-sided 95\% Student-$t$ confidence interval.
A pressure point is marked as converged when the configured minimum number of replicates is available and the relative confidence-interval half-width is below \texttt{convergence.relative\_ci95\_target}.
The outputs are \texttt{convergence\_report.json} and \texttt{convergence\_report.csv}.

Further statistical extensions are:

\begin{itemize}
    \item convergence analysis by plotting loading as a function of simulation step or block index,
    \item block standard error of the mean,
    \item MAE and MSE per reference dataset,
    \item automatic highlighting of the best-matching reference.
\end{itemize}

\subsection{Within-Run Block Convergence}

Completed and still-running GCMC logs can be analyzed without starting another
simulation:

\begin{lstlisting}[language=bash]
python -m analysis.block_convergence \
  outputs/module_A_co2_isotherm/runs/<run_id> \
  --block-size 50000
\end{lstlisting}

The equilibration and production lengths, framework atom count, and adsorbate
size are read from the materialized run. Thermodynamic rows before the end of
equilibration are discarded. The outputs contain block means, cumulative means,
standard errors, progress, and relative changes between the final blocks and
between the half and full production trajectory. Analyses of incomplete logs
are explicitly marked as provisional.

\subsection{Reproducibility Manifest}

Software versions, Git state, system information, the original launch command,
and SHA-256 hashes of immutable run inputs can be archived with:

\begin{lstlisting}[language=bash]
python -m analysis.reproducibility_manifest \
  outputs/module_A_co2_isotherm/runs/<run_id> \
  --config input_json_files/<benchmark.json> \
  --command "python benchmark.py <benchmark.json> ..."
\end{lstlisting}

This writes \texttt{reproducibility/reproducibility\_manifest.json},
\texttt{reproducibility/checksums.sha256}, and, when Conda is available,
\texttt{reproducibility/conda-explicit.txt}. Logs and result files are
intentionally excluded from the immutable-input checksums.

\subsection{Primitive--Conventional Cell Comparison}

Two completed isotherm runs can be compared pressure by pressure with:

\begin{lstlisting}[language=bash]
python -m analysis.cell_comparison \
  outputs/module_A_co2_isotherm/runs/<primitive_run> \
  outputs/module_A_co2_isotherm/runs/<conventional_run> \
  --output-dir reports/cell_comparison
\end{lstlisting}

The report calculates:

\begin{equation}
    R_N =
    \frac{N_{\mathrm{CO_2,conv}}}{N_{\mathrm{CO_2,prim}}},
    \qquad
    R_q =
    \frac{q_{\mathrm{conv}}}{q_{\mathrm{prim}}}.
\end{equation}

The expected molecule-count ratio is inferred from the framework atom counts in
the two LAMMPS data files. It is four for a 424-atom conventional cell compared
with a 106-atom primitive cell. The expected mass-normalized loading ratio is
one. The generated files are \texttt{cell\_comparison.csv},
\texttt{cell\_replicates.csv}, \texttt{cell\_comparison\_summary.json}, and,
when matplotlib is available, \texttt{cell\_comparison.png}. The figure shows
the absolute primitive and conventional molecule-count, mass-normalized
loading, and runtime curves in its upper row. Individual seed replicates are
drawn transparently behind the mean curves. The lower row shows $R_N$, $R_q$,
and the conventional-to-primitive runtime ratio.

New simulations store the wall-clock time of every seed replicate. The cell
comparison reports the mean wall-clock time per replicate and pressure for both
cells, the conventional-to-primitive runtime ratio, and the slowest pressure
points. Runs created before wall-time capture was added can still be compared
for $R_N$ and $R_q$, but their historical runtimes cannot be reconstructed from
the evaluated isotherm alone.
